% ===== PRÉAMBULE CANONIQUE — à coller VERBATIM en tête de CHAQUE .tex =====
\documentclass[11pt,a4paper]{article}
\usepackage[utf8]{inputenc}\usepackage[T1]{fontenc}\usepackage{lmodern}
\usepackage[french]{babel}
\usepackage{amsmath,amssymb,mathtools}\usepackage{icomma}
\usepackage{siunitx}\sisetup{locale=FR}  % \qty{}{}, \si{}, \ang{}
\usepackage{xcolor}\usepackage{colortbl}
\usepackage{tikz}\usepackage{pgfplots}\pgfplotsset{compat=1.18}
\usepackage[siunitx,europeanresistors]{circuitikz} % circuits (élec.)
\usepackage[most]{tcolorbox}\usepackage{enumitem}\usepackage{array,booktabs}
\usepackage[a4paper,margin=2.2cm]{geometry}
\usetikzlibrary{arrows.meta,calc,angles,quotes,positioning,patterns,%
  backgrounds,shapes.geometric,decorations.pathmorphing,decorations.markings,intersections}
\definecolor{primary}{RGB}{222,49,99}\definecolor{primarydark}{RGB}{150,22,55}
\definecolor{defcol}{RGB}{0,114,178}\definecolor{methcol}{RGB}{52,92,148}
\definecolor{excol}{RGB}{0,135,90}\definecolor{attncol}{RGB}{200,40,55}
\tcbset{boxbase/.style={enhanced,breakable,boxrule=0.9pt,arc=2.5pt,
  left=3mm,right=3mm,top=2.3mm,bottom=2.3mm,fonttitle=\bfseries,coltitle=white}}
% --- boîtes TUTORIEL (noms EXACTS, ne pas renommer) ---
\newtcolorbox{cle}[1][]{boxbase,colback=defcol!6,colframe=defcol,
  title={\textsc{À retenir}\ifx\relax#1\relax\relax\else\textmd{ : \itshape #1}\fi}}
\newtcolorbox{methode}[1][]{boxbase,colback=methcol!7,colframe=methcol,sharp corners,
  title={\textsc{Méthode}\ifx\relax#1\relax\relax\else\textmd{ --- \itshape #1}\fi}}
\newtcolorbox{exemplebox}[1][]{boxbase,colback=excol!5,colframe=excol,
  title={\textsc{Exemple}\ifx\relax#1\relax\relax\else\textmd{ : \itshape #1}\fi}}
\newtcolorbox{piege}[1][]{boxbase,colback=attncol!6,colframe=attncol,
  title={\textsc{Piège}\ifx\relax#1\relax\relax\else\textmd{ : \itshape #1}\fi}}
\newtcolorbox{remarque}[1][]{boxbase,colback=black!4,colframe=black!45,
  title={\textsc{Remarque}\ifx\relax#1\relax\relax\else\textmd{ : \itshape #1}\fi}}
% --- boîtes EXERCICE / CORRIGÉ (noms EXACTS, auto-comptées) ---
\newtcolorbox[auto counter]{exo}[1][]{boxbase,colback=primary!4,colframe=primary,
  title={\textsc{Exercice}~\thetcbcounter\ifx\relax#1\relax\relax\else\textmd{ --- \itshape #1}\fi}}
\newtcolorbox[auto counter]{corrige}[1][]{boxbase,colback=excol!4,colframe=excol,
  title={\textsc{Corrigé}~\thetcbcounter\ifx\relax#1\relax\relax\else\textmd{ --- \itshape #1}\fi}}
\newcommand{\vv}[1]{\vec{#1}}\newcommand{\ds}{\displaystyle}
\newcommand{\R}{\mathbb{R}}\newcommand{\N}{\mathbb{N}}\newcommand{\Z}{\mathbb{Z}}
% (un \usepackage{fancyhdr}+en-tête est possible ; il est ignoré côté web)
% =========================================================================
\begin{document}

\begin{center}
{\large\bfseries\color{primarydark} Thème 1 — Niveau Difficile}\\[-2pt]
{\color{primary}\rule{5cm}{1pt}}
\end{center}

\begin{exo}[la longueur d'onde à partir d'un enregistrement]
Une onde sinusoïdale se propage le long d'une corde à la célérité $v=\SI{0.5}{\metre\per\second}$.
On enregistre l'élongation d'\emph{un} point de la corde en fonction du temps.
\begin{center}
\begin{tikzpicture}
\begin{axis}[
  width=11.5cm,height=5.2cm,
  xmin=0,xmax=5,ymin=-12,ymax=12,
  xtick={0,1,2,3,4,5},ytick={-10,0,10},
  grid=both,grid style={black!12},axis lines=left,
  xlabel={$t$ (s)},ylabel={élongation (cm)},
  xlabel style={at={(axis description cs:0.5,-0.13)}},
  tick label style={font=\footnotesize},samples=220,clip=false]
  \addplot[primary,line width=1.3pt,domain=0:5]{10*sin(deg(2*pi*x/2))};
\end{axis}
\end{tikzpicture}
\end{center}
\begin{enumerate}[label=\alph*)]
  \item Lis la période $T$ de l'onde sur le graphe, puis calcule sa fréquence $f$.
  \item Calcule la longueur d'onde $\lambda$ (en \si{\metre} puis en \si{\centi\metre}).
  \item Ce graphe permet-il de lire directement $\lambda$ ? Explique pourquoi.
\end{enumerate}
\end{exo}

\begin{exo}[deux graphes, une vitesse]
Une même onde transversale, se propageant vers les $x>0$, est décrite par \emph{deux} graphes :
à gauche, une photo de la corde à un instant fixé ($y$ en fonction de $x$)~; à droite,
le mouvement d'un point fixe ($y$ en fonction de $t$).
\begin{center}
\begin{tikzpicture}
\begin{axis}[
  width=7.2cm,height=4.7cm,
  xmin=0,xmax=8,ymin=-6,ymax=6,
  xtick={0,2,4,6,8},ytick={-5,0,5},
  grid=both,grid style={black!12},axis lines=left,
  title={\footnotesize photo à $t$ fixé},
  xlabel={$x$ (m)},ylabel={$y$ (cm)},
  xlabel style={at={(axis description cs:0.5,-0.16)}},
  tick label style={font=\footnotesize},samples=200,clip=false]
  \addplot[defcol,line width=1.2pt,domain=0:8]{5*sin(deg(2*pi*x/4))};
\end{axis}
\end{tikzpicture}\hspace{0.5cm}%
\begin{tikzpicture}
\begin{axis}[
  width=7.2cm,height=4.7cm,
  xmin=0,xmax=1.25,ymin=-6,ymax=6,
  xtick={0,0.25,0.5,0.75,1,1.25},ytick={-5,0,5},
  grid=both,grid style={black!12},axis lines=left,
  title={\footnotesize point fixe},
  xlabel={$t$ (s)},ylabel={$y$ (cm)},
  xlabel style={at={(axis description cs:0.5,-0.16)}},
  tick label style={font=\footnotesize},samples=200,clip=false]
  \addplot[primary,line width=1.2pt,domain=0:1.25]{5*sin(deg(2*pi*x/0.5))};
\end{axis}
\end{tikzpicture}
\end{center}
\begin{enumerate}[label=\alph*)]
  \item Sur quel graphe lit-on la longueur d'onde $\lambda$ ? Donne sa valeur.
  \item Sur quel graphe lit-on la période $T$ ? Donne sa valeur.
  \item Calcule la célérité $v$ de l'onde.
\end{enumerate}
\end{exo}

\begin{exo}[échographie : la profondeur d'un obstacle]
Une sonde d'échographie envoie une brève impulsion ultrasonore, qui se réfléchit sur un
organe et revient vers la sonde. L'oscilloscope montre l'impulsion émise (1) et l'écho reçu (2).
La célérité des ultrasons dans les tissus est $v=\SI{1500}{\metre\per\second}$.
\begin{center}
\begin{tikzpicture}[scale=1.0,>=Stealth]
  % axe temps
  \draw[black!70,line width=1pt](0,0)--(10.4,0)node[right,font=\footnotesize]{$t$};
  \foreach \i in {0,...,10}{\draw[black!45](\i,0.1)--(\i,-0.1);}
  % impulsion 1 en division 1
  \draw[primary,line width=1.4pt](1,0)--(1,2.1)--(1.15,0);
  \node[primary,font=\bfseries] at (1,2.4){1};
  % impulsion 2 (écho) en division 9
  \draw[defcol,line width=1.4pt](9,0)--(9,1.3)--(9.15,0);
  \node[defcol,font=\bfseries] at (9,1.6){2};
  % cote entre les deux
  \draw[<->,black](1,-0.6)--(9,-0.6)node[midway,fill=white,inner sep=1pt,font=\footnotesize]{8 divisions};
  \node[below,font=\footnotesize] at (5,-1.05) {1 division $=\SI{10}{\micro\second}$};
\end{tikzpicture}
\end{center}
\begin{enumerate}[label=\alph*)]
  \item Détermine la durée $\Delta t$ qui sépare l'impulsion émise de l'écho.
  \item Cette durée correspond à un trajet \emph{aller-retour}. À quelle profondeur $d$ se
        trouve l'organe ?
\end{enumerate}
\end{exo}

\begin{exo}[rides à la surface de l'eau]
Une pointe frappe la surface de l'eau au point $S$ à rythme régulier et crée des rides
circulaires. On photographie la surface : les cercles en \textbf{trait plein} sont les crêtes,
ceux en \textbf{pointillés} les creux. La longueur d'onde vaut $\lambda=\SI{4}{\centi\metre}$.
Un petit flotteur est placé en $A$.
\begin{center}
\begin{tikzpicture}[scale=1.0]
  \foreach \r in {1,2,3}{\draw[defcol,line width=1pt] (0,0) circle (\r);}
  \foreach \r in {0.5,1.5,2.5}{\draw[defcol!70,dashed,line width=0.7pt] (0,0) circle (\r);}
  \fill[primary] (0,0) circle (2.4pt) node[below left,black,font=\footnotesize]{$S$};
  \coordinate (A) at (35:3);
  \fill[black] (A) circle (2.4pt) node[above right,font=\footnotesize]{$A$};
  \draw[black!55,dashed] (0,0)--(A);
\end{tikzpicture}
\end{center}
Le flotteur $A$ est situé exactement sur la \textbf{3\textsuperscript{e} crête} (le 3\textsuperscript{e}
cercle en trait plein) à partir de $S$.
\begin{enumerate}[label=\alph*)]
  \item En comptant les crêtes, exprime la distance $d=SA$ en fonction de $\lambda$.
  \item Calcule $d$ en centimètres.
\end{enumerate}
\end{exo}

\begin{exo}[un son qui change de milieu]
Un émetteur plongé dans l'eau produit un son de fréquence $f=\SI{250}{\hertz}$. Célérité du son :
\SI{1500}{\metre\per\second} dans l'eau, \SI{340}{\metre\per\second} dans l'air.
\begin{enumerate}[label=\alph*)]
  \item Calcule la longueur d'onde $\lambda_{\text{eau}}$ dans l'eau.
  \item Le son ressort dans l'air. Quelle est alors sa fréquence ? Justifie.
  \item Calcule la longueur d'onde $\lambda_{\text{air}}$ dans l'air, et compare-la à
        $\lambda_{\text{eau}}$ (donne le rapport).
\end{enumerate}
\end{exo}

\end{document}
